\section*{Exercises on \S\arabic{exercisegroup}}
\phantomsection
\addcontentsline{toc}{section}{Exercises on \protect\S\arabic{exercisegroup}}

\begin{enumerate}
\def\labelenumi{\arabic{enumi})}
\tightlist
\item
  Let \(\mathbf{f}\) be a vector function that is \(n\) times differentiable on an interval \(\mathrm{I} \subset \mathbf{R}\) . Prove the formula
\end{enumerate}

\[
\mathrm{D} ^ {n} \mathbf {f} \left(e ^ {x}\right) = \sum_ {m = 1} ^ {n} \frac {a _ {m}}{m !} e ^ {m x} \mathbf {f} ^ {(m)} \left(e ^ {x}\right)
\]

at every point x such that \(e^{x} \in I\) , where the coefficient \(a_{m}\) can be expressed as

\[
a _ {m} = \sum_ {p = 0} ^ {m} (- 1) ^ {p} \binom {m} {p} (m - p) ^ {n}
\]

(method of I, p.~41, exerc. 7, using the Taylor expansion for \(e^{x}\) ).

\begin{enumerate}
\def\labelenumi{\arabic{enumi})}
\setcounter{enumi}{1}
\item
  Let \(f\) be a real function that is \(n\) times differentiable at a point \(x\) , and \(\mathbf{g}\) a vector function \(n\) times differentiable at the point \(f(x)\) . If one puts \(\mathrm{D}^n(\mathbf{g}(f(x))) = \sum_{k=1}^{n} \mathbf{g}^{(k)}(f(x)) u_k(x)\) , then \(u_k\) depends on the function \(f\) alone; deduce from this that \(u_k(x)\) is the coefficient of \(t^k\) in the expansion of \(e^{-tf(x)}\mathrm{D}^n(e^{tf(x)})\) (in terms of \(t\) ).
\item
  For every real \(x > 0\) and every complex \(m = \mu + iv\) one puts \(x^{m} = e^{m\log x}\) ; show that the formula (19) of III, p.~108 remains valid for \(m\) complex and \(x > -1\) and that the remainder \(r_{n}(x)\) in this formula satisfies the inequalities
\end{enumerate}

\[
\begin{array}{l l} | r _ {n} (x) | \leqslant \left| \binom {m - 1} {n} \frac {m}{\mu} x ^ {n} \Big [ (1 + x) ^ {\mu} - 1 \Big ] \right| & \text { if } \mu \neq 0 \\ | r _ {n} (x) | \leqslant \left| \binom {m - 1} {n} m x ^ {n} \log (1 + x) \right| & \text { if } \mu = 0. \end{array}
\]

Generalize the study of the convergence of the binomial series to the case where m is complex.

\begin{enumerate}
\def\labelenumi{\arabic{enumi})}
\setcounter{enumi}{3}
\tightlist
\item
  For every real \(x\) and every number \(p > 1\) prove the inequality
\end{enumerate}

\[
\begin{array}{r l} | 1 + x | ^ {p} \leqslant 1 + p x + \frac {p (p - 1)}{2} x ^ {2} + \dots + \frac {p (p - 1) \dots (p - m + 2)}{(m - 1) !} x ^ {m - 1} \\ & + \frac {p (p - 1) \dots (p - m + 1)}{m !} | x | ^ {m} + h _ {p} | x | ^ {p} \end{array}
\]

where one puts \(m = [p]\) (the integer part of \(p\) ) and

\[
h _ {p} = \frac {p (p - 1) \dots (p - m + 1)}{(m - 1) !} \int_ {0} ^ {1} z ^ {p - m} (1 - z) ^ {m - 1} d z.
\]

¶ 5) Show that \(\pi\) is irrational, in the following way: if one had \(\pi = p/q\) ( \(p\) and \(q\) integers), then on putting \(f(x) = (x(\pi - x))^n / n!\) , the integral \(q^n \int_0^\pi f(x) \sin x dx\) would be an integer \(>0\) (use the formula for integration by parts of order \(n + 1\) ); but show that on the other hand \(q^n \int_0^\pi f(x) \sin x dx\) tends to 0 as \(n\) tends to \(+\infty\) .

\begin{enumerate}
\def\labelenumi{\arabic{enumi})}
\setcounter{enumi}{5}
\tightlist
\item
  Show that on the interval \([-1,+1]\) the function \(|x|\) is the uniform of limit of polynomials, by remarking that \(|x|=(1-(1-x^{2}))^{1/2}\) and using the binomial series. Deduce another proof of the Weierstrass th. from this (cf.~II, p.~83, exerc. 20).
\end{enumerate}

¶ 7) Let \(p\) be a prime number and \(\mathbf{Q}_p\) the field of \(p\) -adic numbers (Gen.~Top., III, p.~322 and 323, exercs. 23 to 25), let \(\mathbf{Z}_p\) be the ring of \(p\) -adic integers, and \(\mathfrak{p}\) the principal ideal \((p)\) in \(\mathbf{Z}_p\) .

\begin{enumerate}
\def\labelenumi{\alph{enumi})}
\tightlist
\item
  Let \(a = 1 + pb\) , where \(b \in \mathbf{Z}_p\) is an element of the multiplicative group \(1 + \mathfrak{p}\) ; show that when the rational integer \(m\) increases indefinitely the \(p\) -adic number \(\frac{(1 + pb)^{p^m} - 1}{p^m}\) tends to a limit equal to the sum of the convergent series
\end{enumerate}

\[
\frac {p b}{1} - \frac {p ^ {2} b ^ {2}}{2} + \dots + (- 1) ^ {n - 1} \frac {p ^ {n} b ^ {n}}{n} + \dots
\]

One denotes this limit by \(\log a\) .

\begin{enumerate}
\def\labelenumi{\alph{enumi})}
\setcounter{enumi}{1}
\item
  Show that when the \(p\) -adic number \(x\) tends to 0 in \(\mathbf{Q}_p\) the number \(\frac{a^x - 1}{x}\) tends to \(\log a\) (use \(a\) ) and the definition of the topology of \(\mathbf{Q}_p\) .
\item
  Show that if \(p \neq 2\) one has \(\log a \equiv pb \pmod{\mathfrak{p}^2}\) , and if \(p = 2\) then \(\log a \equiv 0 \pmod{\mathfrak{p}^2}\) , and \(\log a \equiv -4b^4 \pmod{\mathfrak{p}^3}\) .
\item
  Show that if \(p \neq 2\) (resp. p = 2) then \(x \mapsto \log x\) is an isomorphism of the multiplicative topological group \(1 + p\) onto the additive topological group p (resp. \(p^{2}\) ); in particular, if \(e_{p}\) is the element of \(1 + p\) such that \(\log e_{p} = p\) (resp. \(\log e_{2} = 4\) ) then the isomorphism of \(Z_{p}\) onto \(1 + p\) , which is inverse to \(x \mapsto \frac{1}{p} \log x\) , (resp. \(x \mapsto \frac{1}{4} \log x\) ) is \(y \mapsto e_{p}^{y}\) (cf.~Gen.~Top., III, p.~323, exerc. 25).
\item
  Show that for every \(a \in 1 + p\) the continuous function \(x \mapsto a^{x}\) , defined on \(Z_{p}\) , admits a derivative equal to \(a^{x} \log a\) at every point; deduce from this that the function \(\log x\) admits a derivative equal to 1/x at every point of \(1 + p\) .
\end{enumerate}

¶ 8) a) With the notation of exerc. 7, show that the series with general term \(x^n / n!\) is convergent for every \(x \in \mathfrak{p}\) if \(p \neq 2\) , for every \(x \in \mathfrak{p}^2\) (but for no \(x \notin \mathfrak{p}^2\) ) if \(p = 2\) (determine the exponent of \(p\) in the decomposition of \(n!\) into prime factors). If \(f(x)\) is the sum of this series show that \(f\) is a continuous homomorphism of \(\mathfrak{p}\) (resp. \(\mathfrak{p}^2\) ) into \(1 + \mathfrak{p}^2\) . Deduce from this that for every \(z \in \mathbf{Z}_p\) one has \(f(pz) = e_p^z\) (resp. \(f(p^2z) = e_p^z\) ); in other words,

\[
e _ {p} = 1 + \frac {p}{1 !} + \frac {p ^ {2}}{2 !} + \dots + \frac {p ^ {n}}{n !} + \dots \quad \text { if } p \neq 2
\]

\[
e _ {2} = 1 + \frac {4}{1 !} + \frac {4 ^ {2}}{2 !} + \dots + \frac {4 ^ {n}}{n !} + \dots
\]

(use exerc. 7 e)).

\begin{enumerate}
\def\labelenumi{\alph{enumi})}
\setcounter{enumi}{1}
\tightlist
\item
  For every \(a \in 1 + \mathfrak{p}\) and every \(x \in \mathbf{Z}_p\) show that \(\log (a^x) = x\log a\) , and deduce from a) and exerc. 7 d) that
\end{enumerate}

\[
a ^ {x} = 1 + \frac {x \log a}{1 !} + \frac {x ^ {2} (\log a) ^ {2}}{2 !} + \dots + \frac {x ^ {n} (\log a) ^ {n}}{n !} + \dots
\]

\begin{enumerate}
\def\labelenumi{\alph{enumi})}
\setcounter{enumi}{2}
\tightlist
\item
  For every \(m \in Z_{p}\) show that the continuous function \(x \mapsto x^{m}\) , defined on \(1 + p\) , admits a derivative equal to \(mx^{m-1}\) (use b) and exerc. 7 e)).
\end{enumerate}

\(d\) ) Show that for every \(m\in \mathbf{Z}_p\) and every \(x\in \mathfrak{p}\) if \(p\neq 2\) ( \(x\in \mathfrak{p}^2\) if \(p = 2\) ) the series with general term \(\binom{m}{n}x^n\) is convergent and that its sum is a continuous function of \(m\) ; deduce from this that this sum is equal to \((1 + x)^m\) by remarking that \(\mathbf{Z}\) is (everywhere) dense in \(\mathbf{Z}_p\) .

¶ 9) With the notation of exerc. 7) one denotes by \(\mathbf{O}_2^+ (\mathbf{Q}_p)\) (the group of rotations of the space \(\mathbf{Q}_p^2\) ) the group of matrices of the form

\[
\left( \begin{array}{c c} x & y \\ - y & x \end{array} \right)
\]

with elements in \(\mathbf{Q}_p\) , such that \(x^2 + y^2 = 1\) , this group being endowed with the topology defined in Gen.~Top., VIII, p.~125, exerc. 2.

\begin{enumerate}
\def\labelenumi{\alph{enumi})}
\item
  Denote by \(G_{n}\) the subgroup of \(\mathbf{O}_{2}^{+}(\mathbf{Q}_{p})\) formed by the matrices such that \(t=y/(1+x)\in\mathfrak{p}^{n}\) . Show that \(G_{n}\) is a compact group, that \(G_{n}/G_{n+1}\) is isomorphic to Z/pZ, and that the only compact subgroups of \(G_{1}\) are the groups \(G_{n}\) (cf.~Gen.~Top., III, p.~323, exerc. 24).
\item
  Show that \(G_{1}\) is identical to the subgroup of matrices
\end{enumerate}

\[
\left( \begin{array}{c c} x & y \\ - y & x \end{array} \right)
\]

such that \(x^{2} + y^{2} = 1\) , \(x \in 1 + p^{2}\) and \(y \in p\) if \(p \neq 2\) ; to \(x \in 1 + p^{3}\) and \(y \in p^{2}\) if p = 2.

\begin{enumerate}
\def\labelenumi{\alph{enumi})}
\setcounter{enumi}{2}
\tightlist
\item
  Show that the series with general terms \((-1)^{n} x^{2n}/(2n)!\) and \((-1)^{n-1} x^{2n+1}/(2n+1)!\) are convergent for every \(x \in p\) if \(\neq 2\) , for every \(x \in p^{2}\) if p = 2; let \(\cos x\) and \(\sin x\) be the sums of these series. Show that the map
\end{enumerate}

\[
x \mapsto \left( \begin{array}{c c} \cos x & \sin x \\ - \sin x & \cos x \end{array} \right)
\]

is an isomorphism of the additive topological group \(\mathfrak{p}\) (resp. \(\mathfrak{p}^2\) ) onto the group \(G_1\) .

\(d\) ) If \(p\) is of the form \(4h + 1\) ( \(h\) an integer), there exists in \(\mathbf{Q}_p\) an element \(i\) such that \(i^2 = -1\) . If with every \(z \in \mathbf{Q}_p^*\) one associates the matrix

\[
\left( \begin{array}{c c} \frac {1}{2} \left(z + \frac {1}{z}\right) & \frac {1}{2 i} \left(z - \frac {1}{z}\right) \\ - \frac {1}{2 i} \left(z - \frac {1}{z}\right) & \frac {1}{2} \left(z + \frac {1}{z}\right) \end{array} \right)
\]

one defines an isomorphism of the multiplicative group \(Q_{p}^{*}\) onto the group \(\mathbf{O}_{2}^{+}(\mathbf{Q}_{p})\) ; under this isomorphism the group \(1 + p\) corresponds to \(G_{1}\) and one has \(\cos px + i \sin px = e_{p}^{ix}\) (III, p.~126, exerc. 8).

\begin{enumerate}
\def\labelenumi{\alph{enumi})}
\setcounter{enumi}{4}
\tightlist
\item
  If p is of the form \(4h + 3\) (h an integer) the elements of the matrices of \(\mathbf{O}_{2}^{+}(\mathbf{Q}_{p})\) are necessarily p-adic integers. The polynomial \(X^{2} + 1\) is then irreducible in \(Q_{p}\) ; let \(\mathbf{Q}_{p}(i)\) be the quadratic extension of \(Q_{p}\) obtained by adjoining a root i of \(X^{2} + 1\) ; one endows \(\mathbf{Q}_{p}(i)\) with the topology defined in Gen.~Top., VIII, p.~127, exerc. 2. The group \(\mathbf{O}_{2}^{+}(\mathbf{Q}_{p})\) is isomorphic to the multiplicative group N of elements of \(\mathbf{Q}_{p}(i)\) of norm 1, under the isomorphism which associates the matrix \(\begin{pmatrix} x & y \\ -y & x \end{pmatrix}\) with the element \(z = x + iy\) . Show that there are \(p + 1\) roots of the equation \(x^{p+1} = 1\) in \(\mathbf{Q}_{p}(i)\) and that they form a cyclic subgroup R of N (argue as in Gen.~Top., III, p.~323, exerc. 24; then show that there are \(p + 1\) distinct roots of the congruence \(x^{p+1} \equiv 1 \pmod{p}\) in \(\mathbf{Q}_{p}(i)\) , and, for each root a of this congruence, form the series \((a^{p2n})\) ). Deduce that the group \(\mathbf{O}_{2}^{+}(\mathbf{Q}_{p})\) is isomorphic to the product of the groups R and \(G_{1}\) .
\end{enumerate}

\(f\) ) Show that, for \(p = 2\) , the group \(\mathbf{O}_2^+ (\mathbf{Q}_p)\) is isomorphic to the product of the group \(\mathrm{G}_1\) and the cyclic group of order 4.
% semantic-fragments: legacy-input-seam
\relax{}
